% ============================================================================
% LANGENTWURF LE-08b: El tiempo (Das Wetter)
% ============================================================================
% Fach: Spanisch | Klasse: 7 | Sequenz: 08 (Vacaciones)
% Fachfarbe: #c41e3a (Karminrot)
% Tier: 1 (vollständig)
% ============================================================================

\documentclass[11pt,a4paper]{article}

% ============================================================================
% PACKAGES
% ============================================================================
\usepackage[utf8]{inputenc}
\usepackage[T1]{fontenc}
\usepackage[ngerman]{babel}
\usepackage{geometry}
\usepackage{fancyhdr}
\usepackage{lastpage}
\usepackage{xcolor}
\usepackage{tcolorbox}
\usepackage{enumitem}
\usepackage{tabularx}
\usepackage{longtable}
\usepackage{array}
\usepackage{booktabs}
\usepackage{hyperref}
\usepackage{ifthen}
\usepackage{amssymb}

% ============================================================================
% KONFIGURATION
% ============================================================================

\newcommand{\plantier}{1}
\newcommand{\fachid}{spanisch}
\newcommand{\fach}{Spanisch}
\newcommand{\fachfarbe}{c41e3a}

% ============================================================================
% VARIABLEN
% ============================================================================

\newcommand{\jahrgangsstufe}{7}
\newcommand{\schulart}{Gesamtschule}
\newcommand{\stundenthema}{El tiempo -- Das Wetter}
\newcommand{\reihenthema}{Vacaciones}
\newcommand{\stundennummer}{8b}
\newcommand{\reihenstunden}{10}
\newcommand{\unterrichtsdatum}{XX.XX.XXXX}
\newcommand{\unterrichtszeit}{XX:XX -- XX:XX}
\newcommand{\raum}{XXX}

\newcommand{\lehrkraft}{N.N.}
\newcommand{\ausbildungsstand}{Lehrkraft}

\newcommand{\lerngruppengroesse}{24}
\newcommand{\maennlich}{12}
\newcommand{\weiblich}{12}
\newcommand{\divers}{0}

\newcommand{\gerniveau}{A1}
\newcommand{\lehrwerk}{Qué pasa 1}
\newcommand{\lehrwerkseite}{S. 98--101}

% ============================================================================
% LAYOUT
% ============================================================================
\geometry{left=2.5cm, right=2.5cm, top=2.5cm, bottom=2.5cm}
\definecolor{fach}{HTML}{\fachfarbe}
\definecolor{gray}{HTML}{666666}
\definecolor{lightgray}{HTML}{f5f5f5}
\definecolor{successgreen}{HTML}{2a7a4b}
\definecolor{warningorange}{HTML}{b35c00}

\tcbuselibrary{skins,breakable}

\newtcolorbox{infobox}[1][]{
    colback=lightgray,
    colframe=fach,
    fonttitle=\bfseries,
    title=#1,
    breakable
}

\newtcolorbox{scorebox}{
    colback=white,
    colframe=fach,
    fonttitle=\bfseries,
    title=Didaktik-Score,
    breakable
}

\pagestyle{fancy}
\fancyhf{}
\fancyhead[L]{\small\textcolor{gray}{\fach{} -- Jg. \jahrgangsstufe{} -- \stundenthema}}
\fancyhead[R]{\small\textcolor{gray}{Langentwurf Tier 1}}
\fancyfoot[C]{\small\thepage{} / \pageref{LastPage}}
\renewcommand{\headrulewidth}{0.4pt}
\renewcommand{\footrulewidth}{0pt}

% ============================================================================
% DOKUMENT
% ============================================================================
\begin{document}

% ============================================================================
% DECKBLATT
% ============================================================================
\begin{titlepage}
    \centering
    \vspace*{2cm}

    {\large\textcolor{gray}{\schulart}}\\[2cm]

    {\Huge\textcolor{fach}{\textbf{Langentwurf}}}\\[0.5cm]
    {\large Tier 1 -- Vollständig}\\[2cm]

    {\LARGE\textbf{\stundenthema}}\\[0.5cm]
    {\large im Rahmen der Unterrichtsreihe}\\[0.3cm]
    {\Large\textit{\reihenthema}}\\[2cm]

    \begin{tabular}{rl}
        \textbf{Fach:} & \fach{} (\gerniveau) \\[0.3cm]
        \textbf{Klasse:} & \jahrgangsstufe \\[0.3cm]
        \textbf{Lehrwerk:} & \lehrwerk, \lehrwerkseite \\[0.3cm]
        \textbf{Datum:} & \underline{\hspace{3cm}} \\[0.3cm]
        \textbf{Zeit:} & \underline{\hspace{3cm}} (Doppelstunde) \\[0.3cm]
        \textbf{Raum:} & \underline{\hspace{3cm}} \\[0.3cm]
    \end{tabular}

    \vfill

    {\large Lehrkraft: \lehrkraft}\\[0.3cm]
    {\normalsize\textcolor{gray}{\ausbildungsstand}}\\[1cm]

\end{titlepage}

% ============================================================================
% INHALTSVERZEICHNIS
% ============================================================================
\tableofcontents
\newpage

% ============================================================================
% 1. BEDINGUNGSANALYSE
% ============================================================================
\section{Bedingungsanalyse}

\subsection{Lerngruppe}
\begin{tabular}{ll}
    Kursgröße: & \lerngruppengroesse{} SuS (\maennlich{} m / \weiblich{} w / \divers{} d) \\
    Jahrgangsstufe: & \jahrgangsstufe \\
    Sprachniveau: & \gerniveau{} (GER) \\
    Fremdsprache: & 2. Fremdsprache \\
\end{tabular}

\subsection{Differenzierung (intern)}

\textbf{WICHTIG:} Keine Sternchen auf Schülermaterial!

\begin{tabular}{|l|p{4cm}|p{4cm}|p{4cm}|}
\hline
& \textbf{[B] Basis} & \textbf{[S] Standard} & \textbf{[E] Erweiterung} \\
\hline
\textbf{Ziel} & Berufsreife & Mittlere Reife & Gymnasium \\
\hline
\textbf{Inhalt} & Rezeptiv, Zuordnung & Produktion mit Hilfe & Freie Produktion \\
\hline
\textbf{Hilfen} & Wortschatz, Bilder & Satzanfänge & Nur Impulse \\
\hline
\end{tabular}

\subsection{Besonderheiten}
\begin{itemize}
    \item \textbf{LRS:} 2 SuS (Nachteilsausgleich: Zeitzugabe)
    \item \textbf{DaZ:} 3 SuS (Wortschatzarbeit verstärken, Bildunterstützung)
    \item \textbf{Leistungsstark:} 4 SuS (Zusatzaufgaben, Expertenrolle)
\end{itemize}

% ============================================================================
% 2. SACHANALYSE
% ============================================================================
\section{Sachanalyse}

\subsection{Sprachliche Strukturen}

\textbf{Grammatik -- Wetter mit \textit{hacer}:}
\begin{itemize}
    \item \textbf{hace + sustantivo:} hace sol, hace calor, hace frío, hace viento
    \item \textbf{hace + adjetivo:} hace buen tiempo, hace mal tiempo
    \item Das Verb \textit{hacer} wird unpersönlich verwendet (3. Person Singular)
    \item Keine Übereinstimmung (Genus/Numerus) notwendig
\end{itemize}

\textbf{Grammatik -- Wetter mit \textit{estar}:}
\begin{itemize}
    \item \textbf{estar + adjetivo:} está nublado, está soleado
    \item Bezieht sich auf den Zustand (el cielo, el día)
\end{itemize}

\textbf{Grammatik -- Wetter mit Verben:}
\begin{itemize}
    \item \textbf{llueve} (es regnet) -- von \textit{llover}
    \item \textbf{nieva} (es schneit) -- von \textit{nevar}
    \item Unpersönliche Verben, nur 3. Person Singular
\end{itemize}

\textbf{Wortschatz:}
\begin{itemize}
    \item Wetterausdrücke: hace sol, hace calor, hace frío, hace viento
    \item Wetterverben: llueve, nieva
    \item Zustandsbeschreibung: está nublado
    \item Jahreszeiten: la primavera, el verano, el otoño, el invierno
    \item Fragekonstruktion: ¿Qué tiempo hace?
\end{itemize}

\subsection{Kontrastive Betrachtung (Deutsch $\leftrightarrow$ Spanisch)}
\begin{itemize}
    \item \textbf{Positive Transfers:}
    \begin{itemize}
        \item Unpersönliche Konstruktionen existieren auch im Deutschen (``es regnet'')
        \item Jahreszeiten sind kognat erkennbar (primavera $\sim$ Primär/Frühling)
    \end{itemize}
    \item \textbf{Interferenzen:}
    \begin{itemize}
        \item Deutsch: ``es ist kalt'' vs. Spanisch: ``hace frío'' (wörtlich: ``es macht kalt'')
        \item Kein ``es'' als Subjekt im Spanischen
    \end{itemize}
    \item \textbf{Phonetik:}
    \begin{itemize}
        \item \textit{llueve}: ['llwebe] -- palataler Lateral
        \item \textit{nieva}: ['njeba] -- Diphthong ie
    \end{itemize}
\end{itemize}

\subsection{Kulturelle Aspekte}
\begin{itemize}
    \item Klimazonen Spaniens (Mittelmeer vs. Nordspanien)
    \item Typisches Urlaubswetter an der Costa Brava
    \item Jahreszeiten auf der Südhalbkugel (Lateinamerika)
\end{itemize}

% ============================================================================
% 3. DIDAKTISCHE ANALYSE
% ============================================================================
\section{Didaktische Analyse}

\subsection{Stellung in der Sequenz}
Dies ist die \textbf{8. Doppelstunde (b)} der Sequenz ``\reihenthema'' (\reihenstunden{} Doppelstunden gesamt).

\textbf{Vorwissen:}
\begin{itemize}
    \item Verbkonjugation (ser, estar, tener)
    \item Grundwortschatz Orte und Aktivitäten
    \item Fragen stellen mit ¿Qué...?, ¿Dónde...?
\end{itemize}

\textbf{Anschluss:}
\begin{itemize}
    \item Reiseplanung (08c: Urlaubsaktivitäten)
    \item Dialoge über Ferienpläne
\end{itemize}

\subsection{Kompetenzziele (Rahmenplan MV)}
\begin{itemize}
    \item \textbf{Kommunikativ:} Hörverstehen (Wetterberichte), Sprechen (Wetter beschreiben)
    \item \textbf{Sprachlich:} Unpersönliche Konstruktionen mit hacer, estar
    \item \textbf{Interkulturell:} Klimazonen der spanischsprachigen Welt
    \item \textbf{Methodisch:} Visualisierungen zum Vokabellernen nutzen
\end{itemize}

\subsection{Legitimation}
\textbf{Rahmenplan Spanisch MV (Kl. 7-10):}
\begin{itemize}
    \item Kompetenzbereich: Sprechen (an Gesprächen teilnehmen)
    \item Standard: Einfache Informationen zu vertrauten Themen austauschen
    \item Themenfeld: Alltag, Freizeit, Reisen
\end{itemize}

% ============================================================================
% 4. STUNDENZIELE
% ============================================================================
\section{Stundenziele}

\subsection{Grobziel}
Die SuS erweitern ihre \textbf{kommunikative Kompetenz} im Bereich Sprechen, indem sie das Wetter in verschiedenen Jahreszeiten mit den Konstruktionen \textit{hace + sustantivo}, \textit{estar + adjetivo} und unpersönlichen Verben beschreiben.

\subsection{Feinziele (operationalisiert)}
\begin{tabular}{|c|p{8cm}|c|c|}
\hline
\textbf{Nr.} & \textbf{Die SuS können...} & \textbf{AFB} & \textbf{Diff.} \\
\hline
FZ1 & die Wettervokabeln (hace sol, llueve, etc.) korrekt aussprechen & I & alle \\
\hline
FZ2 & die vier Jahreszeiten auf Spanisch benennen und dem Wetter zuordnen & I & alle \\
\hline
FZ3 & auf die Frage ¿Qué tiempo hace? situationsgerecht antworten & II & alle \\
\hline
FZ4 & einen kurzen Wetterbericht für einen Urlaubsort verfassen & II & [S]/[E] \\
\hline
FZ5 & den Unterschied zwischen hacer und estar beim Wetter erklären & III & [E] \\
\hline
\end{tabular}

% ============================================================================
% 5. METHODISCHE ANALYSE
% ============================================================================
\section{Methodische Analyse}

\subsection{Phasierung (Doppelstunde = 2 × 45 Min.)}
\begin{enumerate}
    \item \textbf{1. Stunde:}
    \begin{itemize}
        \item Einstieg: Bildimpuls (Wetterkarte Spanien)
        \item Erarbeitung: Wettervokabular mit Bildern
        \item Übung: Chorisches Sprechen, Memory-Spiel
        \item Sicherung: Arbeitsblatt Merkkasten
    \end{itemize}
    \item \textbf{2. Stunde:}
    \begin{itemize}
        \item Wiederholung: Wetter-Bingo
        \item Erarbeitung: Jahreszeiten
        \item Anwendung: Mini-Wetterberichte (Partnerarbeit)
        \item Sicherung: Präsentation, AB vervollständigen
    \end{itemize}
\end{enumerate}

\subsection{Medien}
\begin{itemize}
    \item Tafel/Whiteboard
    \item Lehrwerk: \lehrwerk, \lehrwerkseite
    \item Arbeitsblatt: AB-08b.pdf
    \item Lernpfad: LP-08b.html (optional, digital)
    \item Bildkarten: Wettersymbole (16 Karten)
    \item Wetterkarte Spanien (Projektion)
\end{itemize}

% ============================================================================
% 6. VERLAUFSPLANUNG
% ============================================================================
\section{Verlaufsplanung}

\textbf{1. Stunde (45 Min.)}

\begin{longtable}{|p{1cm}|p{1.5cm}|p{4cm}|p{3cm}|p{2cm}|p{2cm}|}
\hline
\textbf{Zeit} & \textbf{Phase} & \textbf{Lehrerhandeln} & \textbf{Schülerhandeln} & \textbf{SF} & \textbf{Medien} \\
\hline
\endhead

0--5 & Einstieg & Zeigt Wetterkarte Spanien: ``Mirad -- ¿qué veis?'' & Vermutungen äußern (deutsch/spanisch) & Plenum & Projektor \\
\hline

5--15 & Input & Präsentiert Wettervokabeln mit Bildern, spricht vor & Hören, Nachsprechen & Plenum & Bildkarten \\
\hline

15--25 & Übung 1 & Leitet chorisches Sprechen an, korrigiert Aussprache & Chorisches/Einzelsprechen & Plenum & -- \\
\hline

25--35 & Übung 2 & Erklärt Memory-Spiel (Bild-Wort), beobachtet & Spielen in 4er-Gruppen & GA & Karten \\
\hline

35--40 & Sicherung & Zeigt AB, erklärt Merkkasten & Tragen Vokabeln ein (Merkkasten) & EA & AB \\
\hline

40--45 & Über\-leitung & Fasst zusammen, Ausblick 2. Stunde & Aufräumen & Plenum & -- \\
\hline
\end{longtable}

\vspace{1em}

\textbf{2. Stunde (45 Min.)}

\begin{longtable}{|p{1cm}|p{1.5cm}|p{4cm}|p{3cm}|p{2cm}|p{2cm}|}
\hline
\textbf{Zeit} & \textbf{Phase} & \textbf{Lehrerhandeln} & \textbf{Schülerhandeln} & \textbf{SF} & \textbf{Medien} \\
\hline
\endhead

0--5 & Wieder\-holung & Wetter-Bingo (ruft Begriffe) & Markieren auf Bingo-Feld & EA & Bingo-Karten \\
\hline

5--15 & Input & Präsentiert Jahreszeiten mit Wetterverbindung & Hören, Notieren & Plenum & Tafel \\
\hline

15--25 & Übung 3 & Zeigt Lernpfad, erklärt Übungen (digital oder AB) & Bearbeiten LP Schritt 3-4 oder AB & EA/PA & Tablets/AB \\
\hline

25--35 & Anwen\-dung & Erklärt Mini-Wetterbericht: ``Estás en Barcelona...'' & Schreiben kurzen Wetterbericht & PA & AB Aufg. 4 \\
\hline

35--40 & Sicherung & Wählt 2-3 Paare zur Präsentation & Präsentieren Wetterbericht & Plenum & -- \\
\hline

40--45 & Abschluss & Fasst Lernziele zusammen, HA & Notieren HA & Plenum & Tafel \\
\hline
\end{longtable}

\textbf{Legende:} EA = Einzelarbeit, PA = Partnerarbeit, GA = Gruppenarbeit, SF = Sozialform

% ============================================================================
% 7. ERWARTETE SCHWIERIGKEITEN
% ============================================================================
\section{Erwartete Schwierigkeiten}

\begin{tabular}{|p{5cm}|p{8cm}|}
\hline
\textbf{Schwierigkeit} & \textbf{Reaktion} \\
\hline
Aussprache \textit{llueve} & Langsam vorsprechen, Zungenposition erklären, chorisches Üben \\
\hline
Verwechslung hace/está & Visualisierung: hace = aktiv (Sonne macht warm), está = Zustand (Himmel ist bewölkt) \\
\hline
Deutsche Interferenz ``es ist kalt'' & Bewusstmachen: ``hacer'' = ``machen'' $\rightarrow$ ``es macht kalt'' \\
\hline
Zeitproblem & Puffer: Memory-Spiel kürzen, Präsentationen reduzieren \\
\hline
Heterogenität & [B] nur Zuordnung, [S] Lückentext, [E] freie Produktion \\
\hline
\end{tabular}

% ============================================================================
% 8. HAUSAUFGABE
% ============================================================================
\section{Hausaufgabe}

\begin{itemize}
    \item Lehrwerk S. 109, Aufgabe 3 (Wetterdialog ergänzen)
    \item Vokabeln lernen: hace sol bis invierno
    \item Optional: Lernpfad LP-08b.html vollständig bearbeiten
\end{itemize}

% ============================================================================
% 9. ANLAGEN
% ============================================================================
\section{Anlagen}

\begin{enumerate}
    \item Arbeitsblatt (AB-08b.pdf)
    \item Musterlösung (ML-08b.pdf)
    \item Lehrerhinweise (LH-08b.pdf)
    \item Lernpfad (LP-08b.html)
    \item Bildkarten Wetter (16 Stück)
    \item Bingo-Karten (Klassensatz)
    \item Wetterkarte Spanien (Projektion)
\end{enumerate}

% ============================================================================
% 10. DIDAKTIK-SCORE
% ============================================================================
\newpage
\section{Didaktik-Score}

\begin{scorebox}
\textbf{Bewertung der didaktischen Qualität nach 10 Dimensionen}\\
\textbf{Ziel-Score: $\geq$ 9.0 (Premium-Qualität)}

\vspace{1em}
\begin{tabular}{|c|l|c|c|p{4.5cm}|}
\hline
\textbf{\#} & \textbf{Dimension} & \textbf{Gew.} & \textbf{Score} & \textbf{Notiz} \\
\hline
1 & Differenzierung & 15\% & 9/10 & [B]/[S]/[E] qualitativ verschieden \\
\hline
2 & Taxonomie (AFB) & 10\% & 9/10 & I: 40\%, II: 40\%, III: 20\% \\
\hline
3 & Lernziel-Alignment & 10\% & 10/10 & Rahmenplan MV Kl. 7 passend \\
\hline
4 & Aktivierung & 10\% & 9/10 & Memory, Bingo, Partnerarbeit \\
\hline
5 & Scaffolding & 10\% & 9/10 & Bilder, Wortlisten, Satzanfänge \\
\hline
6 & Feedback-Qualität & 10\% & 9/10 & Chorisches Üben, Peer-Feedback \\
\hline
7 & Sprachliche Klarheit & 10\% & 10/10 & Altersgerecht, klare Anweisungen \\
\hline
8 & Motivation \& Relevanz & 10\% & 9/10 & Urlaubskontext, Wetterkarte \\
\hline
9 & Inklusion (UDL) & 10\% & 9/10 & Bilder, Audio, Bewegung \\
\hline
10 & Praktikabilität & 5\% & 9/10 & 90 Min. realistisch, Material vorhanden \\
\hline
\multicolumn{3}{|r|}{\textbf{GESAMT:}} & \textbf{9.2/10} & \textcolor{successgreen}{\textbf{FREIGABE}} \\
\hline
\end{tabular}

\vspace{1em}
\textbf{Bewertungsschwellen:}
\begin{itemize}
    \item[$\boxtimes$] \textcolor{successgreen}{$\geq$ 9.0: \textbf{FREIGABE} (Premium-Qualität)}
    \item[$\Box$] \textcolor{warningorange}{8.0--8.9: Iteration empfohlen}
    \item[$\Box$] 6.0--7.9: Überarbeitung erforderlich
    \item[$\Box$] $<$ 6.0: Grundlegende Neukonzeption
\end{itemize}
\end{scorebox}

\end{document}
